
\section{Background and Related Work}
\label{background}

SECDA (SystemC Enabled Co-Design of DNN Accelerators) is a hardware-software co-design methodology aimed at reducing the development time of FPGA-based accelerators for DNN inference~\cite{haris2021secda}. It provides accelerator design, system simulation, and application framework integration within a unified development environment. 

SECDA-TFLite~\cite{haris2023secda} extends the SECDA methodology by providing a toolkit for integrating FPGA accelerators into the TensorFlow Lite (TFLite) framework directly. It includes modules for SystemC simulation, profiling, and AXI-based communication, enabling developers to rapidly prototype and evaluate new accelerator designs. Through the use of the TFLite delegate system, SECDA-TFLite allows portions of a DNN model to be offloaded to custom hardware accelerators while maintaining compatibility with the existing inference pipeline. 

Despite these advances, the process of identifying efficient accelerator designs remains challenging and time consuming even for experts, since it involves exploring numerous architectural parameters such as compute unit dimensions, memory organization, tiling strategies, and dataflow patterns. Manually exploring these parameters needs multiple iterations and often relies on domain expertise~\cite{dhilleswararao2022efficient, boutros2025field}. 
However, the recent progress in Large Language Models (LLMs) offers new opportunities to automate parts of the DSE process without the need for extensive domain expertise~\cite{ling2025domain, zheng2025automation}. 
For example, iDSE~\cite{li2025idse} proposes an LLM-based framework for navigating HLS design spaces, where the model guides the search process and prunes candidate configurations to converge more efficiently toward optimal solutions. 
Similarly, LUMINA~\cite{zhang2026luminallmguidedgpuarchitecture} proposed an LLM-driven architecture exploration framework for GPU design, where the model analyzes simulator code and performs bottleneck-based reasoning. 
Also, LIMCA~\cite{vungarala2025limca} proposes an LLM-driven framework to explore analog in-memory computing architectures, where the model generates and evaluates configurations. 

While these works demonstrate the growing potential of LLMs in assisting hardware architecture exploration across domains, SECDA-DSE differs by (i) integrating LLM reasoning directly within the accelerator design workflow and (ii) enabling iterative DSE guided by a hardware evaluation feedback loop.